This section is partially formalized. The new companion uses Mathlib's actual complex exponential \texttt{Function.Periodic.qParam 1}, whose defining formula is the book's $e(z)=\exp(2\pi iz)$. Its restriction to the actual upper half-plane represents $q(z)$. No group quotient, cusp width, length or area is defined to equal the intended conclusion. Steven's correspondence review remains pending.

The \hyperref[modular.cusp-parameter]{cusp parameter} has checked defining formula, nonvanishing, radius, integer periodicity, and invariance under the actual modular action of $T^n$. The fiber identity \eqref{eq:cusp-parameter-fibers} is expressed as existence of an integer $n$ with $z=w+n$. The radius condition is equivalent to the height cutoff, and every nonzero point of the prescribed disk has a preimage in the upper half-plane. Surjectivity is proved for all $Y\geq0$; the modular-neighborhood embedding in the prose requires the stronger $Y>1$ and remains open. The limit as imaginary part tends to infinity lands at the puncture $0$.

The shared library proves the fiber and preimage results for any nonzero or positive real translation period, as appropriate. It reuses Mathlib's logarithmic inverse and its inverse-modulo-period theorem. The prose instead separates equality of radii and angles and constructs an inverse from a locally chosen argument. The formal inverse uses the principal complex logarithm and is only used to obtain a preimage; no assertion that it is a continuous global inverse is made. The full trigonometric polar expression is not yet a companion theorem, although its radius and fiber consequences are checked.

Remaining targets are the exact elliptic and parabolic stabilizers, rational parabolic fixed points and the unique cusp orbit, cusp width as a property of an effective stabilizer with a specified coordinate normalization, the height estimate excluding other modular identifications, continuity and the quotient homeomorphism, the horizontal cross-section length and cusp area, and noncompactness. The horizontal-segment length formula \eqref{eq:horizontal-segment-length} is a prose calculation; its path-length formalization remains open in Section 1.3.

The \hyperref[modular.quotient-topology]{quotient-topology theorem} additionally requires the cusp compactification topology, the side-pairing model, the cut into two closed disks, and the gluing homeomorphism to the sphere. The two cone points remain in the underlying topological space; only the cusp is removed. These topological statements are not yet kernel-checked.

Machine evidence is recorded in \texttt{lean/audits/chapter01-release-check-2026-11-12/record.json}. Consult its outcome and exact source hashes. This is partial section coverage, not a completed human correspondence review.
